\scanpage{122}
\unsure{correspondent|correspondant} exactement (via $x\longmapsto\underline{O}_{\mathfrak X,x}$) aux jauges dans $L$ qui sont ent. de type fini sur $K$, de corps des fractions $L$ --- i.e aux $V$ précédents.

Sauf erreur, ces deux types de p$^t$ de $\mathfrak X$ (qui épuisent \deleted{tout} si $n=1$, \unsure{)} \inserted{qui est le seul cas, avec $n=0$, ou $\mathfrak X$ soit un \unsure{schém.|schéma} !)} sont les seuls pour lesquels $\underline{O}_{\mathfrak X,x}$ soit noethérien. Cependant, on a

\margin{$\big|$}
\uline{Proposition} Pour $\forall x\in\mathfrak X$, $\underline{O}_{\mathfrak X,x}$ est un anneau de valuation dans $L$, de corps des fractions $L$, et contenant $K$. L'application $x\longmapsto\underline{O}_{\mathfrak X,x}$ est une bijection de $\mathfrak X$ sur l'ens de ces anneaux dans $L$.
\margin{On pourrait partir d'un schéma excellent $X$, et construire un $\mathfrak X$ relatif au dessus\ldots}

\uline{Dém} a) Soit $f\in L\setminus\underline{O}_{\mathfrak X,x}$, il faut prouver $1/f\in\underline{O}_{\mathfrak X,x}$. Or $f=g/h$, \inserted{$g,h$}$\unsure{\in}\underline{O}_{X,a}$, où $X$ modèle \inserted{\unsure{régulier|regulier}} (propre), $a$ un p$^t$, $f,g$ \deleted{\unsure{premiers} entre eux} --- utilisant la résolution, \deleted{\unsure{on} \uline{sépare} \unsure{doit|deux} les \unsure{diviseurs}} \inserted{OPS $\exists$ système régulier} de paramètres, avec $f,g$ des monômes
\margin{121}
